\pdfoutput=1
\pdfmapfile{+lm-ec.map}
\pdfmapfile{+lm-math.map}
\pdfmapfile{+lm-ts1.map}
\documentclass[10pt,letterpaper]{article}
\usepackage[margin=0.62in,top=0.53in,bottom=0.53in]{geometry}
\usepackage[T1]{fontenc}
\usepackage{lmodern,microtype,enumitem,titlesec,hyperref,fancyhdr}
\hypersetup{colorlinks=true,urlcolor=black,pdfauthor={Zixian Liu},pdftitle={Zixian Liu — Curriculum Vitae}}
\setlength{\parindent}{0pt}
\setlength{\parskip}{0pt}
\renewcommand{\baselinestretch}{1.04}
\titleformat{\section}{\large\bfseries\scshape}{}{0pt}{}[\vspace{-4pt}\rule{\linewidth}{0.35pt}]
\titlespacing*{\section}{0pt}{10pt}{5pt}
\setlist[itemize]{leftmargin=12pt,label=\textbullet,itemsep=2pt,topsep=3pt,parsep=0pt,partopsep=0pt}
\pagestyle{fancy}\fancyhf{}\renewcommand{\headrulewidth}{0pt}
\fancyfoot[R]{\small Zixian Liu \enspace | \enspace \thepage}
\newcommand{\entry}[2]{\textbf{#1}\hfill #2\par}
\newcommand{\subentry}[2]{\vspace{4pt}\textbf{#1}\hfill #2\par}
\begin{document}
\begin{center}
{\LARGE\bfseries ZIXIAN LIU}\\[4pt]
\href{mailto:liu3866@purdue.edu}{liu3866@purdue.edu}\quad|\quad\href{https://kevin-liu-pu.github.io}{kevin-liu-pu.github.io}
\end{center}\vspace{-7pt}
\section{Education}
\entry{Purdue University}{West Lafayette, IN}
B.S. in Computer Science; Minor in Mathematics\hfill Sep. 2023--May 2027 (expected)\\
GPA: 3.94/4.0; Dean's List and Semester Honors in five semesters.
\begin{itemize}
\item Selected coursework: Data Structures, Systems Programming, Linear Algebra, Statistics; Fall 2026: Data Mining and Machine Learning, Analysis of Algorithms, Operating Systems, Numerical Methods.
\end{itemize}
\section{Manuscripts Under Review}
[1] First-author manuscript on real-world robot learning and manipulation, under review at IEEE ICRA 2027. (Title omitted during double-anonymous review.)
\section{Research Experience}
\entry{Machine Learning and Robotics Research}{Jan. 2026--present}
Purdue University\hfill Advisor: Prof. Yung-Hsiang Lu
\begin{itemize}
\item Team lead since Summer 2026; coordinate five undergraduates across two research tracks in experiment planning and results analysis; participate in weekly graduate research meetings.
\end{itemize}
\subentry{Multimodal Robot Learning and Cross-String Adaptation}{Aug. 2026--present}
\begin{itemize}
\item Developing a policy architecture that integrates multimodal observations, scores from a supervised human-preference model, and musical score requirements to generate bowing commands for a UR5e robotic cellist.
\item Investigating reward design and evaluation using audio, force/torque, vision, and robot state, with independent listening studies to assess performance.
\item Extending the bowing action space beyond speed and depth to bow angle and contact position; conducting experiments on transfer across strings and adaptation to string-specific contact conditions.
\end{itemize}
\subentry{Vision-Based Human--Multi-Robot Collaboration}{May 2026--present}
\begin{itemize}
\item Developed a MediaPipe/OpenCV prototype to track hand movements and detect beats from downward gestures; built a metronome that updates tempo and click timing from detected beats in real time.
\item Tested tempo changes, pauses, and tracking interruptions; developing a vision interface for a human conductor to coordinate a robotic cello and glockenspiel.
\end{itemize}
\subentry{Residual RL for Precise, Contact-Rich Manipulation [1]}{Jan. 2026--Sep. 2026}
\begin{itemize}
\item Developed a four-stage residual SAC framework for a UR5e robotic cellist under contact drift, using a fixed reward classifier trained on human ratings of the robot's recordings.
\item Designed held-out classifier evaluation and a blinded paired-listening study; policies achieved a 15\% mean held-out classifier-score gain over the planner across five pieces.
\item Built the real-robot training stack with RTDE control, audio capture, contact calibration, logging, force/torque-based limits, and watchdog restarts; wrote the first-author manuscript and a reproducibility audit.
\end{itemize}
\subentry{Thermal and Force Signals for Material and Tactile Property Recognition}{Spring 2025}
Dalian University of Technology\hfill Advisor: Prof. Bo Zhang
\begin{itemize}
\item Investigated whether thermal signals add predictive information beyond force and other nonthermal tactile signals across CLAMP, PHAC-2, and SENS3 public datasets.
\item Identified group-level differences hidden by pooled CLAMP results; examined class support, sensor baselines, thermal changes, and normalization, then proposed complementary datasets to test task-dependent effects.
\item Evaluated a three-class SENS3 pilot using 38 contact records: adding temperature change and heat flux increased macro recall from 38.3\% with force alone to 95.1\% on held-out surfaces.
\item Tested stricter cross-participant transfer, obtaining 57.5\% and 69.4\% macro recall; distinguished held-out-surface performance from participant generalization and noted the confounding reduction in training data.
\end{itemize}
\newpage
\section{Selected Projects}
\entry{Xinu Operating System}{Aug. 2026--present}
\begin{itemize}
\item Implemented semaphore-based synchronization and independent interprocess communication with bounded message buffers and blocking send/receive operations; deployed and tested on ARM BeagleBone Black.
\item Built a kernel-level semaphore checker to identify inconsistencies among semaphore counts, wait queues, and process states.
\end{itemize}
\subentry{Unix Shell}{Jan. 2026--May 2026}
\begin{itemize}
\item Built a Unix shell in C/C++ with lex/yacc, supporting pipelines, input/output redirection, background jobs, subshells, and signal handling.
\item Implemented variable and wildcard expansion and command history.
\end{itemize}
\subentry{Web Crawler and Search Engine}{Aug. 2025--May 2026}
\begin{itemize}
\item Built a crawler and search engine with indexing, hash-table and AVL-tree data structures, and multi-keyword queries.
\item Benchmarked query performance across data structures and maintained the system beyond the course project.
\end{itemize}
\subentry{Systematic Trend-Timing Strategy}{Feb. 2026--Sep. 2026}
\begin{itemize}
\item Built a vectorized backtester over 100 years of daily data with next-open execution, transaction costs, and taxes.
\item Evaluated 233 variants using preregistered walk-forward tests, block bootstrap, and multiple-testing correction; identified selection overfitting and reduced crisis drawdowns without excess returns.
\end{itemize}
\section{Industry Experience}
\entry{Guosen Securities Company Ltd.}{Shenzhen, China}
Quantitative Analyst Intern\hfill Jun. 2022--Sep. 2022
\begin{itemize}
\item Developed a Dynamic Conditional Correlation (DCC) model in Python and R to project forward correlations of Russell 2000 stocks.
\item Built Tableau and Power BI reports over more than three million daily data entries for the management team and external investors.
\end{itemize}
\subentry{Xiaowo Technology Co. Ltd. (a China Unicom subsidiary)}{Shanghai, China}
Software Engineer, Wobao Department\hfill Mar. 2021--Dec. 2021
\begin{itemize}
\item Integrated SQL workflows with Python for daily data-warehouse updates and cross-department reporting.
\item Implemented event-driven pipelines for high-frequency client data and optimized Django back-end modules for the client web portal.
\end{itemize}
\subentry{Huamei Frontier Logistics Company Ltd.}{Dalian, China}
Data Science Intern\hfill Dec. 2020--Mar. 2021
\begin{itemize}
\item Improved internal logistics systems with PL/SQL pre-QA checks and validation; estimated the revenue impact of new distribution hubs using difference-in-differences analysis.
\end{itemize}
\section{Technical Skills}
\textbf{Programming:} Python, C/C++, Java, R, MATLAB, SQL, \LaTeX.\\[4pt]
\textbf{Robot learning and ML:} PyTorch, SAC, PPO, residual reinforcement learning, learned reward models, audio feature extraction (librosa), held-out evaluation, blinded listening studies.\\[4pt]
\textbf{Robotics and perception:} UR5e / ur\_rtde, force/torque sensing, contact calibration, safety supervisors, MediaPipe, OpenCV.\\[4pt]
\textbf{Systems and tools:} Linux, Git, lex/yacc, ffmpeg, Django, Tableau, Power BI.
\end{document}
